\documentclass[10pt]{article}
\usepackage[margin=0.60in]{geometry}
\usepackage{amsmath,booktabs,enumitem,fancyhdr,graphicx,listings,tabularx,titlesec,xcolor}
\usepackage[hidelinks]{hyperref}
\graphicspath{{docs/lessons/assets/}{assets/}}
\hypersetup{
 pdftitle={ADK Lesson 065: Qualified Sets of Identified 18B20 Probes},
 pdfauthor={Kevin Shanahan},
 pdfsubject={Qualifying copied search, conversion, and scratchpad evidence for four configured DS18B20 identities},
 pdflang={en-US}}
\pagestyle{fancy}\fancyhf{}
\lhead{\textsc{ADK Lesson 065}}\rhead{Qualified 18B20 Probe Sets}
\cfoot{\thepage}\setlength{\headheight}{14pt}
\setlength{\parindent}{0pt}\setlength{\parskip}{3pt}
\setlist{nosep,leftmargin=1.45em}
\titleformat{\section}{\Large\bfseries}{\thesection}{0.5em}{}
\newcommand{\lessonbox}[1]{\begin{center}\fbox{\parbox{0.92\linewidth}{#1}}\end{center}}
\newcommand{\blankline}{\rule{0.76in}{0.4pt}}
\newcommand{\longblank}{\rule{1.35in}{0.4pt}}
\lstset{
 basicstyle=\ttfamily\footnotesize,
 columns=fullflexible,
 breaklines=true,
 frame=single,
 rulecolor=\color{black},
 showstringspaces=false}

\begin{document}
\small
\begin{center}
 {\Huge\bfseries Qualified Sets of Identified 18B20 Probes}\\[3pt]
 {\Large Keep four identities stable while evidence changes}\\[7pt]
 \fbox{\textbf{E0 SYNTHETIC REPLAY --- NO BUS, PROBE, OR THERMAL CLAIM}}
\end{center}
\begin{tabularx}{\linewidth}{@{}>{\bfseries}l X >{\bfseries}l X@{}}
Lesson & 065 --- component & Time & 180--240 minutes\\
Prerequisite & Lesson 064 copied transactions & Board & compile-only Mega replay\\
Input & copied search/conversion/read evidence & Observation & named snapshot cells\\
\end{tabularx}

\lessonbox{\textbf{Responsibility.}
\texttt{Qualified18B20ProbeSetPolicy} qualifies one bounded copied cycle
against exactly four configured family-\texttt{0x28} ROM identities. It
publishes four stable attributed slots with signed temperature code-space
intervals, age, quality, and status. It does not discover, power, sample, or
calibrate a physical sensor.}

\section{Qualify evidence, not a box of parts}
Lesson 064 proves that a bounded policy can correlate copied single-wire
transaction evidence. Lesson 065 gives some of those transactions DS18B20
meaning: a chained ROM search, a conversion start and completion observation,
and a nine-byte scratchpad read. Meaning arrives only after structure,
identity, chronology, and CRC checks agree.

\begin{center}
% ADK visual: pencil
\includegraphics[width=0.96\linewidth,height=1.95in,keepaspectratio]
 {065-evidence-funnel-pencil.png}

\small\textit{Alternative text: four graphite measuring bars with progressively
finer tick marks sit beside four closed intervals that narrow from a broad
shaded span to a single narrow band, illustrating resolution without claiming
physical accuracy.}
\end{center}

\begin{tabularx}{\linewidth}{@{}X X@{}}
\toprule
E0 can establish & E0 cannot establish\\
\midrule
family byte and Dallas ROM CRC & genuine package, maker, or cable identity\\
copied scratchpad CRC and resolution bits & electrical timing, rise time, or supply current\\
stable configured-slot attribution & physical position or probe response time\\
deterministic freshness and step policy & thermal accuracy, waterproofing, or safe use\\
\bottomrule
\end{tabularx}

There is no wiring diagram or formal schematic at E0. A future powered lesson
must qualify exact specimens, pin order, external supply, ordinary pull-up,
rails, line timing, four-device loading, and independent non-Serial evidence.
Parasite power, strong pull-up, immersion, medical, food, preservation, and
unattended heat-source claims remain excluded.

\newpage
\section{Configure exactly four stable identities}
Configuration contains four distinct complete ROM codes. Each must use family
byte \texttt{0x28} and pass the Dallas CRC-8 over its first seven bytes. The
algorithm begins at zero, processes each byte least-significant bit first,
uses reflected polynomial \texttt{0x8c}, and applies no final XOR.
\footnote{\href{https://www.analog.com/media/en/technical-documentation/data-sheets/ds18b20.pdf}{Analog Devices, \emph{DS18B20 Programmable Resolution 1-Wire Digital Thermometer}, Rev. 6.}}

\begin{center}
% ADK visual: pencil
\includegraphics[width=0.96\linewidth,height=1.85in,keepaspectratio]
 {065-four-identity-slots-pencil.png}

\small\textit{Alternative text: four numbered graphite storage bins each hold
a deliberately different token shaped like a small probe, a visual metaphor
for four identities that retain their assigned slots rather than depictions of
qualified physical specimens.}
\end{center}

Each slot also configures 9-, 10-, 11-, or 12-bit resolution, a conversion
deadline, maximum age, inclusive signed raw-sixteenth range, and maximum
per-cycle step magnitude. There is no runtime count, empty wildcard, or
replacement identity. Search order cannot change configured order.

\begin{tabularx}{\linewidth}{@{}l l l X@{}}
\toprule
Resolution & Undefined raw bits & Maximum conversion & Represented interval\\
\midrule
9 bit  & 3 & 94 ms  & \([m,m+7]\) sixteenths\\
10 bit & 2 & 188 ms & \([m,m+3]\) sixteenths\\
11 bit & 1 & 375 ms & \([m,m+1]\) sixteenths\\
12 bit & 0 & 750 ms & \([m,m]\) sixteenths\\
\bottomrule
\end{tabularx}

Here \(m\) is the signed raw value after masking undefined low bits. The closed
interval describes digital code space, not sensor accuracy or uncertainty.
The final API requires the configured deadline to equal the resolution maximum
for its slot; a shorter or longer value rejects configuration.

\begin{enumerate}
 \item A 9-bit scratchpad decodes to masked raw value \(400\). Record the
       interval: \blankline\ through \blankline.
 \item A search returns the same four ROMs in reverse order. Which configured
       slot changes identity? \blankline.
 \item Why is an all-zero optional ROM incompatible with this component?
       \longblank.
\end{enumerate}

\section{Stage one complete cycle}
A caller lifetime-reuses one opaque \texttt{Ds18b20CycleBuilder}. After
\texttt{beginCycle()}, role-specific ingestion methods accept one complete
Lesson 064 snapshot at a time. \texttt{finishSearch()} records whether
enumeration completed or explicitly exceeded the four-result evidence
capacity. Only \texttt{finalizeCycle()} may commit policy state.

\begin{center}
% ADK visual: pencil
\includegraphics[width=0.96\linewidth,height=1.65in,keepaspectratio]
 {065-staged-cycle-pencil.png}

\small\textit{Alternative text: a graphite line runs from a loose probe-shaped
token through a bounded geared mechanism, an hourglass, and a paper roll,
representing copied input, staged processing, supplied time, and one published
result rather than a physical circuit.}
\end{center}

\newpage
\section{Make search continuity prove presence}
The builder retains at most four
\texttt{Ds18b20NormalizedSearchPass} records. Each record contains the
normalized transaction reference plus its requested and completed
\texttt{OneWireSearchState}.
The first request uses empty search state. Every later request must equal the
immediately preceding completed result, including ROM,
\texttt{lastDiscrepancy}, and \texttt{lastDevice}. No pass may follow a
terminal result. The request/result chain is the bounded policy rendering of
the manufacturer's discrepancy-tree search procedure.\footnote{\href{https://www.analog.com/en/resources/app-notes/1wire-search-algorithm.html}{Analog Devices, \emph{1-Wire Search Algorithm}, Application Note 187.}}

\begin{center}
% ADK visual: pencil
\includegraphics[width=0.96\linewidth,height=1.95in,keepaspectratio]
 {065-search-chain-pencil.png}

\small\textit{Alternative text: eight distinct graphite beads form one
connected chain, ending in a bead with a check mark; a magnifying glass above
suggests inspecting each linked search result and its terminal evidence.}
\end{center}

A complete nonempty enumeration ends with \texttt{lastDevice} true. One
through four returned devices may therefore prove configured identities
missing. A zero-pass cycle cannot prove all four missing because it contains
no completed search witness. Four results ending nonterminal may explicitly
prove over-capacity; the fifth device is never silently dropped.

\begin{tabularx}{\linewidth}{@{}l X X@{}}
\toprule
Search evidence & Set result & What it does not imply\\
\midrule
complete terminal chain, known distinct ROMs & \texttt{Complete} or
\texttt{Missing} & physical location\\
duplicate returned ROM & \texttt{DuplicateIdentity} & two interchangeable probes\\
CRC-valid foreign ROM & \texttt{UnknownIdentity} & permission to replace a slot\\
bad-family or bad-CRC returned ROM & \texttt{TransportFault} & attributable configured identity\\
truncated, failed, or over-capacity & \texttt{TransportFault} & proved absence\\
\bottomrule
\end{tabularx}

\texttt{presentMask} marks configured identities proved present by a
structurally complete search. \texttt{validCount} counts only
\texttt{Current} slots. \texttt{faultMask} marks every slot that is not
\texttt{Current}. These fields answer different questions and must not be
collapsed.

Predict a complete search that finds three configured ROMs and omits the
fourth: set quality \blankline; missing slot still has its configured ROM?
\blankline; array length \blankline.

\section{Bind conversion and read evidence}
For each returned device, conversion start, conversion-status continuation,
and scratchpad read must agree on exact ROM, owner, lifecycle, configuration,
predecessor transaction, and one nonzero conversion generation. A crossed
witness rejects the entire cycle atomically.

\begin{center}
% ADK visual: pencil
\includegraphics[width=0.96\linewidth,height=1.55in,keepaspectratio]
 {065-correlation-braid-pencil.png}

\small\textit{Alternative text: a graphite clock precedes four pinned evidence
notes along one path; the path rises to two final notes, one clean and one
shaded, illustrating that chronology and correlation determine the admitted
outcome.}
\end{center}

Completion requires a correlated conversion-status receipt whose copied bit
is high. Elapsed time alone does not prove completion. A high bit may arrive
before the configured maximum. A low bit exactly at or after that maximum, or
a high bit after it, is transport-fault evidence.

\newpage
\section{Decode scratchpads without inventing precision}
A DS18B20 scratchpad is exactly nine bytes. Byte eight must equal the Dallas
CRC-8 of bytes zero through seven before temperature or resolution is used.
Configuration byte bit 7 must be zero and bits 4--0 must be one; bits 6--5
encode resolution and must match the configured slot.

\begin{center}
% ADK visual: pencil
\includegraphics[width=0.96\linewidth,height=1.95in,keepaspectratio]
 {065-scratchpad-interval-pencil.png}

\small\textit{Alternative text: a pencil points to a graphite notebook strip
of ordered byte-like cells whose final cell bears an integrity shield and
check mark, emphasizing validation before interpretation.}
\end{center}

The published numeric unit remains a signed sixteenth of a degree Celsius.
Negative values use the signed 16-bit scratchpad representation. Widened
arithmetic forms intervals and absolute step differences so endpoints do not
overflow.

\begin{tabularx}{\linewidth}{@{}l X@{}}
\toprule
Quality & Interpretation\\
\midrule
\texttt{Current} & correlated completion and CRC-valid, configured evidence\\
\texttt{ScratchpadCrcFault} & nine bytes arrived, but integrity failed\\
\texttt{ResolutionMismatch} & scratchpad resolution disagrees with its slot\\
\texttt{ResetDefaultWithoutConversion} & raw \(+85^\circ\mathrm{C}\) lacks exact completed conversion\\
\texttt{ImplausibleStep} & widened absolute step exceeds configured maximum\\
\texttt{Stale} & last trusted value is older than maximum age\\
\texttt{Missing} & complete search proves configured identity absent\\
\texttt{TransportFault} & correlated transport/chronology evidence failed\\
\bottomrule
\end{tabularx}

Raw \(+85^\circ\mathrm{C}\) is not automatically bad. Without exact correlated
completion it is the reset-default quality; after correlated completion it is
a possible current reading. Likewise, an implausible step is a domain warning,
not proof that a physical event was impossible. The new valid sample becomes
the next comparison baseline.

\begin{enumerate}
 \item One scratchpad byte changes but its old CRC is retained. Which numeric
       fields may refresh? \blankline.
 \item At the maximum configured step, is the new value admitted?
       \blankline. One sixteenth beyond? \blankline.
 \item Does a code-space interval authorize an accuracy claim? \blankline.
\end{enumerate}

\newpage
\section{Preserve trusted values through pending and missing states}
Freshness begins at the caller-supplied policy-clock time passed with the
correlated scratchpad. Lesson 065 also retains transaction-clock times, but
does not invent an epoch mapping between the two clocks. Age equal to
\texttt{maximumAge} remains current; one millisecond older is stale.

\begin{center}
% ADK visual: pencil
\includegraphics[width=0.96\linewidth,height=1.95in,keepaspectratio]
 {065-quality-history-pencil.png}

\small\textit{Alternative text: four numbered positions support separate
stacks of graphite tags on pegs; the stacks differ in depth while positions
one through four remain fixed, illustrating stable identity with retained
history.}
\end{center}

\texttt{ConversionPending} retains the last trusted raw value, interval,
observation time, and increasing age. It does not read an incomplete
scratchpad or refresh freshness. \texttt{Missing} likewise retains the stable
identity and prior trusted value. A later fully correlated current cycle
returns that same identity to its original slot.

Set-level structural and transport decisions precede per-slot qualities. A
valid cycle commits all four slots together, so one CRC fault does not erase
three current readings. Status precedence follows configured slot order after
set-level failures; status and quality remain separate facts.

\begin{tabularx}{\linewidth}{@{}l l l l@{}}
\toprule
Slot & Search & Evidence & Published quality\\
\midrule
0 & present & complete, CRC valid & \texttt{Current}\\
1 & present & completion still low, before limit & \texttt{ConversionPending}\\
2 & absent & complete terminal search & \texttt{Missing}\\
3 & present & scratchpad CRC bad & \texttt{ScratchpadCrcFault}\\
\bottomrule
\end{tabularx}

For this table, write \texttt{validCount} \blankline\ and the four-bit
\texttt{faultMask} \blankline. Which rows may still retain an older trusted
numeric value? \longblank.

\section{Reject structural changes atomically}
Cycle sequence is nonzero and modularly forward; \texttt{UINT32\_MAX} followed
by 1 is its defined nonzero rollover. Zero, regression, backward time, and an
exact half-range ambiguity reject. An equal sequence accepts only a
field-for-field identical completed builder image.

A fieldwise-identical normalized duplicate is idempotent: it cannot refresh age, extend a
pending conversion, advance step history, or change disappearance evidence.
A changed duplicate rejects without modifying policy state or the caller's
output. Ingestion failure leaves every byte of the staged builder unchanged;
failed finalization leaves policy and output unchanged.

\newpage
\section{Replay the policy on a Mega without touching hardware}
The canonical sketch creates four fixed configurations, replays copied
Lesson 064 snapshots into one reusable builder, finalizes one cycle, and
copies selected results into named volatile cells. Compile and inspect those
cells; do not upload this E0 exercise to claim a powered fixture.

\begin{center}
% ADK visual: pencil

\includegraphics[width=0.96\linewidth,height=1.85in,keepaspectratio]
 {065-mega-replay-pencil.png}

\small\textit{Alternative text: an open graphite notebook has four numbered
rows, distinct tool-like identity symbols, blank result rectangles, and
different geometric quality marks; a pencil suggests recording the replay
outcome without depicting a wired board.}
\end{center}

\begin{enumerate}
 \item \textbf{Predict:} record configured slot order, expected qualities,
       \texttt{presentMask}, \texttt{faultMask}, and \texttt{validCount}.
 \item \textbf{Observe:} compile for Mega 2560, then inspect the named volatile
       result cells in a debugger, simulator, or disassembly. Serial is
       optional supporting context, not the only evidence.
 \item \textbf{Interpret:} compare each stable ROM slot, quality, raw interval,
       age, and status with the copied cycle. Explain any disagreement before
       changing fixture data.
\end{enumerate}

This observation establishes software acquisition and deterministic result
publication only. There is no electrical safe-state experiment because the E0
component owns no endpoint. Future E1 resource acquisition and physical
released-line evidence are separate checks.

\section{Publisher verification and resource boundary}
\begin{tabularx}{\linewidth}{@{}>{\bfseries}l X@{}}
Mega replay & ordinary: 13,662 B flash, 1,438 B static SRAM; isolated no-LTO:
16,196 B flash, 1,438 B static SRAM.\\
Objects and stack & policy 764 B; reusable builder 477 B; snapshot 180 B;
conservative synchronous stack 533 B.\\
Placement & recurring policy plus builder 1,241 B; lifetime placement including
one snapshot 1,421 B; residual SRAM 6,093 B after static, stack, and reserve.\\
Evidence identity & \texttt{350ddd1f6cefef2f5da249f2956323391}\newline
\texttt{0faf9fb11e9c874c9b0377c768ef0bc3}\\
Publication evidence & Publisher-owned PDF validation and rendered pencil
visual review passed; independent resource rereview reproduced the evidence
cleanly.\\
Physical evidence & Deferred; no powered specimen or bus is authorized.\\
\end{tabularx}

Planning hard limits are 16 KiB ordinary replay flash, 1,536 B static SRAM,
640 B conservative synchronous stack, 768 B policy object, 512 B reusable
builder, and 512 B caller-owned snapshot/result. Every hard gate and the
residual floor pass. Ordinary flash and SRAM, stack, policy, and builder miss
their planning targets but remain below their hard limits; independent review
accepted those bounded misses without weakening replay identity or recurring
operation. Exact specimen, electrical, and physical E1 acceptance remain open.

\section{Exercises and next boundary}
\begin{enumerate}
 \item Construct a complete search permutation that leaves slot order
       unchanged. Mark which evidence proves one configured ROM missing.
 \item Corrupt each scratchpad byte in turn. Confirm that no corruption
       partially refreshes value, age, presence, or step baseline.
 \item Replay pending immediately before, exactly at, and immediately after
       each resolution deadline. Explain the boundary outcomes.
 \item Demonstrate current, pending, stale, missing, and reappeared states for
       one ROM without ever changing its configured index.
 \item Explain why Lesson 066 may use configured ROMs for spatial order but
       cannot use search order as physical position.
\end{enumerate}

Lesson 066 will consume the complete copied four-slot snapshot and map a
configured subset into spatial order, adjacent code-space gradients,
fault-dominant presentation pages, and caller-owned record intent. It still
will not drive a display, read an RTC, write storage, or turn these digital
intervals into calibrated thermal accuracy.

\textit{Stop at E0: qualify exact powered specimens and the complete
four-device bus before adding wiring, schematics, or physical claims.}
\end{document}
